\section{Introduction}
\label{sec:intro}
% claim: I01
A motion monitor for a repetitive robot must account for where the robot is in its task. The same joint position can be expected at one stage and inconsistent with another. Existing robot-monitoring studies use temporal alignment, prediction or execution progress to make this distinction: CASPER evaluates motion anomalies using external inertial measurements, while progress-dependent monitoring and Bayesian interaction primitives explicitly represent temporal state~\cite{casper2024,park2016,campbell2017}. These precedents motivate a precise question about an observer's output: does the evidence contained in a measurement correction improve detection beyond its prospective prediction error?

% claim: I02
Prediction and correction serve different purposes. A predictor estimates the next observation from available history. A measurement update then changes the state estimate to account for the new observation. Innovation-based failure detection and the relationship between tracking sensitivity and residuals are established topics~\cite{willsky1976}. Likewise, posterior-to-prior divergence is an established measure of Bayesian surprise~\cite{itti2005}. Their existence does not imply that combining a residual with a correction statistic will improve a particular security monitor. A normal observation can legitimately change a phase estimate, and an abnormal observation can be poorly explained without greatly redistributing phase weights.

% claim: I03
We examine that question on the original CASPER right-arm recording using joint positions, joint velocities and relative elapsed time. A fixed empirical phase observer supplies a prospective residual and a Jensen--Shannon correction discrepancy. The primary comparison adds the discrepancy to the residual while holding the observer trajectory, observations, fitted reference and evaluation support constant. Both variants already incorporate nominal progression. This comparison therefore isolates an added score term, not the presence of progression or a new state estimator.

% claim: I04
The study provides three bounded outputs: an explicit same-observer comparison with classical surprise and ordinary controls; a complete chronological evaluation that separates normal calibration from final outcomes; and a reproducible record of the resulting detection--false-alarm trade-offs. All eleven deterministic methods are retained. The scope is a descriptive study of one acquisition, not a claim that phase estimation or score fusion is new, nor a reproduction of published CASPER performance. We use the term \emph{empirical motion observer}: the implementation has measured-state correspondence and causal synchronization, but no validated robot dynamics or independently calibrated physical-state uncertainty.

% claim: I05
The findings support only a small correction-score effect. Absolute detection is weak, and simple controls match or exceed several observer outcomes. Reporting this result requires more than counting event hits. A detector that alarms often during normal motion can intersect every long intervention interval while recognizing few of its positive observations. We therefore examine point metrics, alarm episodes, complete event outcomes and schedule-relative delays together. This choice follows the broader need to align security-learning evaluation with the information and operational setting being claimed~\cite{arp2022}.
